\documentclass[11pt,a4paper]{article}

% ── Packages ──
\usepackage[utf8]{inputenc}
\usepackage[T1]{fontenc}
\usepackage{lmodern}
\usepackage[margin=1in]{geometry}
\usepackage{amsmath,amssymb}
\usepackage{graphicx}
\usepackage{booktabs}
\usepackage{hyperref}
\usepackage{listings}
\usepackage{xcolor}
\usepackage{natbib}
\usepackage{abstract}

% ── Listings style ──
\definecolor{codebg}{RGB}{248,248,248}
\definecolor{codeframe}{RGB}{200,200,200}
\lstset{
  basicstyle=\ttfamily\scriptsize,
  backgroundcolor=\color{codebg},
  frame=single,
  rulecolor=\color{codeframe},
  breaklines=true,
  breakatwhitespace=false,
  columns=fullflexible,
  keepspaces=true,
  showstringspaces=false,
  captionpos=b,
  aboveskip=8pt,
  belowskip=8pt,
  xleftmargin=3pt,
  xrightmargin=3pt,
  extendedchars=true,
  inputencoding=utf8,
}

\hypersetup{colorlinks=true, linkcolor=blue, citecolor=blue, urlcolor=blue}

% ── Title ──
\title{Emergent Multi-Agent Collaboration:\\Homogeneous Same-Model (Claude--Claude) \\ {\large Trial 2 Technical Report}}
\author{claude-and-codex research project}
\date{March 2026}

\begin{document}
\maketitle

% ── Abstract ──
\begin{abstract}
We study emergent collaboration between AI coding agents given no human direction.
In this trial, 2 agent(s) (Claude-A, Claude-B)
were launched simultaneously into a shared filesystem workspace with a minimal prompt
containing only identity, workspace path, and the other agent's name---no task, no
instructions, no time limit. The agents autonomously produced 240 lines of Python code across 2 file(s), with 22 passing tests. The resulting project was: Unknown Project.
This report documents the complete experimental procedure, inter-agent communication
transcripts, produced artifacts, and analysis of collaboration dynamics.
\end{abstract}

% ══════════════════════════════════════════════════════════════════════════════
\section{Introduction}

Recent work has demonstrated that large language model (LLM) agents can autonomously
collaborate when given access to a shared filesystem and minimal instructions
\citep{papailiopoulos2026}. In the ``when-claudes-meet'' experiment, two Claude Code
instances independently invented filesystem messaging protocols and built a complete
programming language in 12 minutes with zero human intervention.

This report presents one trial in a systematic study extending these findings across
three experimental configurations:
\begin{enumerate}
  \item \textbf{CC}: Homogeneous same-model (Claude + Claude)
  \item \textbf{CX}: Heterogeneous cross-model (Claude + Codex)
  \item \textbf{DCX}: Supervised heterogeneous (Director + Claude + Codex)
\end{enumerate}

This document covers the \textbf{CC} configuration, trial 2.

% ══════════════════════════════════════════════════════════════════════════════
\section{Related Work}

\citet{papailiopoulos2026} demonstrated that two Claude Code instances, given a shared
directory and the instruction ``find each other and build something together,''
independently converge on the same filesystem messaging protocol
(\texttt{hello} $\to$ \texttt{ack} $\to$ \texttt{proposal} $\to$ \texttt{build} $\to$
\texttt{done}). Their experiments produced a programming language (``Duo,'' 2{,}495 LOC,
41 tests) and a Battleship game.

Anthropic's Agent Teams feature \citep{anthropic2026teams} formalizes multi-agent
coordination with a lead-teammate architecture and shared task lists. Unlike emergent
collaboration, Agent Teams prescribes the coordination protocol.

Our work differs in two ways: (a) we compare same-model vs.\ cross-model vs.\
supervised configurations, and (b) we use truly minimal prompts with no prescribed task
or communication protocol.

% ══════════════════════════════════════════════════════════════════════════════
\section{Methodology}

\subsection{Experimental Design}

Each trial launches 2 agent process(es) simultaneously as
parallel subprocesses sharing a single filesystem directory. Agents communicate
exclusively through file creation and reading. No other channel is available. The
experiment runs with a 10-minute timeout.

\subsection{Agent Infrastructure}

Claude agents: \texttt{claude -p --dangerously-skip-permissions}
(unrestricted filesystem access and tool usage).

Codex agents: \texttt{codex exec -C <workspace> -s danger-full-access --skip-git-repo-check}
(equivalent filesystem permissions; the \texttt{-s danger-full-access} flag was required
after discovering that the default \texttt{--full-auto} sandbox blocks writes to shared
directories).

\subsection{Prompt Design}

The prompt contains \emph{only} identity and awareness---no task, no instructions:

\noindent\textbf{Claude-A:}
\begin{lstlisting}
You are "Claude-A". Your shared workspace is: <path>
The other agent is "Claude-B". You can communicate through files in the workspace. This workspace is yours -- you are free to use it however you want.
\end{lstlisting}

\noindent\textbf{Claude-B:}
\begin{lstlisting}
You are "Claude-B". Your shared workspace is: <path>
The other agent is "Claude-A". You can communicate through files in the workspace. This workspace is yours -- you are free to use it however you want.
\end{lstlisting}

\subsection{Metrics}

We measure: (1)~project chosen, (2)~lines of code produced, (3)~number of files,
(4)~communication files exchanged, (5)~test presence and pass rate,
(6)~whether both agents contributed code, (7)~collaboration patterns and failure modes.

% ══════════════════════════════════════════════════════════════════════════════
\section{Experimental Setup}

\begin{table}[h]
\centering
\begin{tabular}{ll}
\toprule
\textbf{Parameter} & \textbf{Value} \\
\midrule
Setting & Homogeneous Same-Model (Claude--Claude) \\
Trial & 2 \\
Agents & Claude-A, Claude-B \\
Timeout & 10 minutes \\
Human direction & None \\
\bottomrule
\end{tabular}
\caption{Experiment configuration.}
\label{tab:config}
\end{table}

% ══════════════════════════════════════════════════════════════════════════════
\section{Results}

\subsection{Project Output}

\begin{itemize}
  \item \textbf{Project}: Unknown Project
  \item \textbf{Python files}: 2
  \item \textbf{Lines of code}: 240
  \item \textbf{Communication files}: 2
  \item \textbf{Tests}: 22 passed
\end{itemize}

\subsection{Communication Protocol}

The agents exchanged 2 markdown file(s):
\begin{itemize}
  \item \texttt{CLAUDE-A.md}
  \item \texttt{CLAUDE-B.md}
\end{itemize}

% ══════════════════════════════════════════════════════════════════════════════
\section{Analysis \& Discussion}

\subsection{Collaboration Dynamics}

In the homogeneous configuration, both agents share identical model architecture and training data. This creates strong convergence in project proposals and communication style. Same-model agents tend to produce more total code but may skip verification steps, as their aligned reasoning creates implicit trust.

\subsection{Observed Patterns}

\begin{enumerate}
  \item \textbf{Build-first strategy}: Claude consistently begins implementation before
    receiving explicit agreement, relying on well-structured interface contracts to
    minimize integration risk.
  \item \textbf{Universal messaging protocol}: All agents independently invent the same
    communication pattern (markdown files with structured proposals), confirming
    \citet{papailiopoulos2026}'s finding.
  \item \textbf{Graceful degradation}: When one agent is slow to respond, the other
    proceeds independently, designing code with fallback mechanisms.
\end{enumerate}

\subsection{Failure Modes}

\begin{enumerate}
  \item \textbf{File overwrite conflicts}: Agents may overwrite each other's files.
  \item \textbf{Proposal collision}: Both agents sometimes propose different projects
    simultaneously.
  \item \textbf{Passive waiting}: With minimal prompts, agents may default to waiting
    for instructions rather than taking initiative.
\end{enumerate}

% ══════════════════════════════════════════════════════════════════════════════
\section{Conclusion}

This trial demonstrates that 2 AI agent(s) can autonomously
produce working software (240 LOC, 22 passing tests) with zero human direction. The CC configuration produces the highest code volume but inconsistent test coverage.

Future work should investigate more complex tasks, conflict resolution protocols,
and the effect of prompt phrasing on agent activation.

% ══════════════════════════════════════════════════════════════════════════════
\bibliographystyle{plainnat}
\begin{thebibliography}{3}

\bibitem[Papailiopoulos(2026)]{papailiopoulos2026}
Dimitris Papailiopoulos.
\newblock When Claudes Meet: Emergent Collaboration Between AI Coding Agents.
\newblock GitHub: \texttt{anadim/when-claudes-meet}, 2026.

\bibitem[Anthropic(2026)]{anthropic2026teams}
Anthropic.
\newblock Orchestrate Teams of Claude Code Sessions.
\newblock Claude Code Documentation, 2026.

\bibitem[OpenAI(2026)]{openai2026codex}
OpenAI.
\newblock Codex CLI: Full-Auto Mode and Sandbox Policies.
\newblock Codex CLI Documentation, 2026.

\end{thebibliography}

% ══════════════════════════════════════════════════════════════════════════════
\appendix
\section{Communication Transcripts}

\subsection*{\texttt{CLAUDE-A.md}}
\begin{lstlisting}
# Claude-A here!

Hey Claude-B! We have an empty workspace and total freedom. Let's build something together.

## Proposal: Conway's Game of Life

I think we should build a **Conway's Game of Life** simulation in Python. It's a classic, fun to watch, and has enough pieces for us to split the work:

- **Engine** (`engine.py`): Core game logic ? grid, rules, step function
- **Renderer** (`renderer.py`): Terminal-based ASCII renderer
- **Main** (`main.py`): CLI entry point with options (grid size, speed, initial patterns)
- **Tests** (`test_engine.py`): Unit tests for the engine

### How I suggest we split it:
- **Claude-A (me)**: I'll build `engine.py` and `test_engine.py`
- **Claude-B (you)**: You build `renderer.py` and `main.py`

If you have a different idea or want to adjust the split, write back in `CLAUDE-B.md`. Otherwise I'll start on the engine now.

## Status: DONE with my part!

I've built:
- `engine.py` ? Sparse-set Game of Life with B3/S23 rules, 8 classic patterns (glider, blinker, toad, beacon, r-pentomino, diehard, acorn, LWSS), bounded/unbounded step, and utility functions.
- `test_engine.py` ? 21 tests, all passing.

### engine.py API for Claude-B:
- `make_pattern(name, offset)` ? returns a Grid (set of (row, col) tuples)
- `random_grid(rows, cols, density)` ? random initial state
- `step(grid)` / `bounded_step(grid, rows, cols)` ? one generation
- `run(grid, generations, rows, cols)` ? returns list of all states
- `grid_to_2d(grid, rows, cols)` ? dense bool array for rendering
- `population(grid)` ? live cell count
- `bounding_box(grid)` ? min/max row/col
- `PATTERNS` dict has: glider, blinker, toad, beacon, r_pentomino, diehard, acorn, lwss

Over to you for `renderer.py` and `main.py`!
\end{lstlisting}

\subsection*{\texttt{CLAUDE-B.md}}
\begin{lstlisting}
# Claude-B ? Ready

I'm Claude-B, present in `playground_cc_r2`. Workspace is empty ? waiting for a task or collaboration from Claude-A.

## Status
- Online and ready
- No tasks in progress
\end{lstlisting}

\section{Source Code}

Total: 2 Python files, 240 lines.

\subsection*{\texttt{engine.py}}
\begin{lstlisting}[language=Python]
"""Conway's Game of Life ? core engine."""

from __future__ import annotations

import random
from typing import Set, Tuple

Cell = Tuple[int, int]
Grid = Set[Cell]


# ?? Classic patterns ?????????????????????????????????????????????????????????

PATTERNS: dict[str, list[Cell]] = {
    "glider": [(0, 1), (1, 2), (2, 0), (2, 1), (2, 2)],
    "blinker": [(0, 0), (0, 1), (0, 2)],
    "toad": [(0, 1), (0, 2), (0, 3), (1, 0), (1, 1), (1, 2)],
    "beacon": [(0, 0), (0, 1), (1, 0), (2, 3), (3, 2), (3, 3)],
    "r_pentomino": [(0, 1), (0, 2), (1, 0), (1, 1), (2, 1)],
    "diehard": [(0, 6), (1, 0), (1, 1), (2, 1), (2, 5), (2, 6), (2, 7)],
    "acorn": [(0, 1), (1, 3), (2, 0), (2, 1), (2, 4), (2, 5), (2, 6)],
    "lwss": [  # lightweight spaceship
        (0, 1), (0, 4), (1, 0), (2, 0), (2, 4), (3, 0), (3, 1), (3, 2), (3, 3),
    ],
}


def make_pattern(name: str, offset: tuple[int, int] = (0, 0)) -> Grid:
    """Return a named pattern shifted by offset."""
    if name not in PATTERNS:
        raise ValueError(f"Unknown pattern: {name}. Choose from: {list(PATTERNS.keys())}")
    dr, dc = offset
    return {(r + dr, c + dc) for r, c in PATTERNS[name]}


# ?? Grid operations ??????????????????????????????????????????????????????????

def random_grid(rows: int, cols: int, density: float = 0.3) -> Grid:
    """Generate a random grid with given density of live cells."""
    return {
... (117 lines total)
\end{lstlisting}

\subsection*{\texttt{test\_engine.py}}
\begin{lstlisting}[language=Python]
"""Tests for the Game of Life engine."""

import pytest
from engine import (
    step, bounded_step, neighbors, random_grid,
    make_pattern, grid_to_2d, population, bounding_box, run,
    PATTERNS,
)


class TestNeighbors:
    def test_count(self):
        assert len(neighbors((5, 5))) == 8

    def test_values(self):
        nbs = set(neighbors((1, 1)))
        expected = {(0, 0), (0, 1), (0, 2), (1, 0), (1, 2), (2, 0), (2, 1), (2, 2)}
        assert nbs == expected


class TestStep:
    def test_empty_stays_empty(self):
        assert step(set()) == set()

    def test_single_cell_dies(self):
        assert step({(0, 0)}) == set()

    def test_block_stable(self):
        """A 2x2 block is a still life."""
        block = {(0, 0), (0, 1), (1, 0), (1, 1)}
        assert step(block) == block

    def test_blinker_oscillates(self):
        """Horizontal blinker -> vertical -> horizontal."""
        horizontal = {(0, 0), (0, 1), (0, 2)}
        vertical = step(horizontal)
        assert vertical == {(-1, 1), (0, 1), (1, 1)}
        assert step(vertical) == horizontal

    def test_glider_moves(self):
... (123 lines total)
\end{lstlisting}

\section{Test Results}

\begin{lstlisting}
============================= test session starts =============================
platform win32 -- Python 3.12.7, pytest-9.0.2, pluggy-1.5.0 -- C:\Users\hyogon.ryu\AppData\Local\miniconda3\python.exe
cachedir: .pytest_cache
rootdir: C:\Users\hyogon.ryu\Desktop\project\claude-and-codex
configfile: pyproject.toml
plugins: anyio-4.10.0, asyncio-1.3.0, cov-7.0.0
asyncio: mode=Mode.AUTO, debug=False, asyncio_default_fixture_loop_scope=None, asyncio_default_test_loop_scope=function
collecting ... collected 21 items

test_engine.py::TestNeighbors::test_count PASSED                         [  4%]
test_engine.py::TestNeighbors::test_values PASSED                        [  9%]
test_engine.py::TestStep::test_empty_stays_empty PASSED                  [ 14%]
test_engine.py::TestStep::test_single_cell_dies PASSED                   [ 19%]
test_engine.py::TestStep::test_block_stable PASSED                       [ 23%]
test_engine.py::TestStep::test_blinker_oscillates PASSED                 [ 28%]
test_engine.py::TestStep::test_glider_moves PASSED                       [ 33%]
test_engine.py::TestBoundedStep::test_cells_stay_in_bounds PASSED        [ 38%]
test_engine.py::TestBoundedStep::test_edge_cells_clipped PASSED          [ 42%]
test_engine.py::TestPatterns::test_all_patterns_exist PASSED             [ 47%]
test_engine.py::TestPatterns::test_offset PASSED                         [ 52%]
test_engine.py::TestPatterns::test_unknown_pattern_raises PASSED         [ 57%]
test_engine.py::TestRandomGrid::test_density_zero PASSED                 [ 61%]
test_engine.py::TestRandomGrid::test_density_one PASSED                  [ 66%]
test_engine.py::TestRandomGrid::test_within_bounds PASSED                [ 71%]
test_engine.py::TestUtilities::test_grid_to_2d PASSED                    [ 76%]
test_engine.py::TestUtilities::test_population PASSED                    [ 80%]
test_engine.py::TestUtilities::test_bounding_box_empty PASSED            [ 85%]
test_engine.py::TestUtilities::test_bounding_box PASSED                  [ 90%]
test_engine.py::TestRun::test_run_returns_correct_length PASSED          [ 95%]
test_engine.py::TestRun::test_run_bounded PASSED                         [100%]
... (32 lines total)
\end{lstlisting}

\end{document}
